\section{Discussion}\label{sec:discussion}

\paragraph{The representation is not the bottleneck.} The oracle condition shows that a small
ontology, rendered into a few hundred tokens, carries nearly everything the probes ask for, that
every engine we tested can read it, and that, unlike a transcript, it does not degrade as the
history grows. What limits
ontology-mediated migration with small models is write-in: deciding which utterances carry lasting
information, whom they are about, and whether they replace something already stored. This is
encouraging for the architecture, because extraction is the part that improves with every model
generation and can be run by a different, stronger engine than the one the user talks to, while
the ontology and the store stay fixed. Our strong-extractor runs confirm it: a store written once
by a hosted model is read by 7--14B engines as well as the gold store. This suggests a deployment
in which a capable model curates memory at write time, and any engine, including a small local
one, carries it.

\paragraph{When the transcript is enough.} A frontier reader recovers nearly everything from a
long transcript, so for a user who only ever moves between frontier engines the ontology's case is
not accuracy but cost and control: \res{ptok-ratio-a2x-b-large}$\times$ fewer prompt tokens per
question, provenance per item, and erasure that removes a value instead of adding a request to
forget it next to the value.

\paragraph{Alignment loss.} A shared ontology avoids pairwise alignment, but only for what it can
express. Anything outside \memcore's classes is either forced into a \texttt{Fact} with a free-text
attribute, where it is portable but weakly typed, or dropped. Our enumerations are closed (eight
relation types); a user's grandparent or neighbour does not fit. Extending the ontology is a
SemVer-minor change that the registry checks, but every adapter must then move to the new version.
LLM-based matching~\cite{qiang2024agentom} could map vendor memory schemas onto \memcore where a
vendor keeps its own.

\paragraph{Governance.} Who maintains the shared ontology matters more than its first version.
Publishing it as a versioned package with competency questions as tests, content hashes and a
persistent namespace gives the governance problem a concrete form: changes are proposed as
versions, their compatibility is checked mechanically, and consumers pin what they rely on.
A standards body, such as the W3C Community Group on agent memory~\cite{w3cmemcg2026}, is a
natural owner.

\paragraph{Vendor incentives and native export.} Vendors gain from lock-in and have little reason
to adopt a shared memory format unprompted. Native memory exports, where they exist, are free-text
lists or proprietary JSON, closer to our summary condition than to an ontology store; we could not
measure them with local engines. We expect regulation to be the adoption mechanism: portability
rights~\cite{gdpr,quebec25} already require a machine-readable export, and an ontology-backed export
is the cheapest way for a vendor to meet that requirement in a form other engines can actually use.

\paragraph{Erasure.} The tombstone design makes erasure auditable, but our results show that
delegating the decision to the engine is risky in both directions: missed requests leak, and
over-broad ones destroy memory. Erasure belongs in the store's API, confirmed by the user, with the
engine at most proposing candidates.
